% คำนำ — หัวข้อ "คำนำ" และชื่อผู้จัดทำท้ายหน้าสร้างโดย environment preface ใน main.tex
% ตามแนวทางของ template: ระบบอะไร รุ่นใด ผู้อ่านแต่ละกลุ่มควรอ่านบทใด และช่องทางแจ้งข้อผิดพลาด (ไม่เกิน 1 หน้า)

เอกสารคู่มือการใช้งานฉบับนี้จัดทำขึ้นเพื่ออธิบายการติดตั้ง การใช้งาน การดูแลรักษา และการพัฒนาต่อยอดโปรเจกต์ ไดโน: เว็บแอปพลิเคชันเพื่อการท่องเที่ยวจังหวัดขอนแก่น ซึ่งเป็นระบบที่ช่วยให้นักท่องเที่ยวสามารถค้นหาสถานที่และกิจกรรมท่องเที่ยว สอบถามข้อมูลผ่านแชตบอต วางแผนการเดินทาง และสะสมพอยท์ผ่านการสแกน QR Code ขณะเดียวกันยังช่วยให้เจ้าหน้าที่สามารถจัดการและปรับปรุงข้อมูลการท่องเที่ยวให้มีความถูกต้องและเป็นปัจจุบัน

คู่มือฉบับนี้แบ่งเนื้อหาออกเป็นบทตามกลุ่มผู้อ่าน เพื่อให้สามารถศึกษาและนำระบบไปใช้งานได้อย่างเหมาะสม โดยนักท่องเที่ยวควรศึกษา บทที่ 3 การเริ่มต้นใช้งาน และบทที่ 4 การใช้งานสำหรับนักท่องเที่ยว ส่วนเจ้าหน้าที่ผู้รับผิดชอบการจัดการข้อมูลผ่านหน้าเว็บควรศึกษา บทที่ 3 และบทที่ 5 การใช้งานสำหรับผู้ดูแลระบบ สำหรับผู้ที่รับผิดชอบการติดตั้งและดูแลเครื่องแม่ข่ายควรศึกษา บทที่ 2 ความต้องการของระบบและการติดตั้ง และบทที่ 6 การดูแลระบบ ขณะที่นักพัฒนาที่รับช่วงต่อควรศึกษา บทที่ 2 บทที่ 7 คู่มือสำหรับนักพัฒนา และภาคผนวก ข ทั้งนี้ ผู้อ่านทุกกลุ่มควรเริ่มต้นจากบทที่ 1 ภาพรวมของระบบ เพื่อทำความเข้าใจวัตถุประสงค์และความสามารถของระบบ และสามารถศึกษาเพิ่มเติมได้จากบทที่ 8 เมื่อพบปัญหาหรือข้อขัดข้องในการใช้งาน

ผู้จัดทำหวังเป็นอย่างยิ่งว่าคู่มือฉบับนี้จะเป็นประโยชน์ต่อผู้ใช้งานทุกกลุ่ม ช่วยให้สามารถใช้งานและดูแลระบบได้อย่างถูกต้อง มีประสิทธิภาพ และสะดวกยิ่งขึ้น